\section{Mass and Chirality}

A fermion has mass. An electron weighs something, and that weight is the reason matter is heavy and the reason atoms hold their size. In most treatments the mass is a number you write down. You declare it, you call it the electron mass, and the theory carries it along without ever saying where it came from. That is a hole. A substrate that wants to earn its physics must produce the mass as a mechanism, not hold a slot for a value.

This section closes that hole on the substrate $\{3,4,3,4\}$. The claim is sharp. The fermion mass is the coupling between the two chiralities of the dock's sites, the $8s$ and the $8c$. It is read off the substrate two independent ways that agree to machine precision, and it vanishes into a massless Weyl fermion the moment the coupling is switched off.

\begin{principle}[Mass is a rate of trade]
The mass is not an assumed energy. It is the rate at which the two chiralities trade. Switch the trade off and the mass is exactly zero.
\end{principle}

\subsection{The hole the mass used to be}

It pays to name what was wrong before, because the shape of the old error tells you what a real fix must supply. In the earlier tests the fermion mass was simply assigned. A number was chosen and inserted, and the rest of the calculation ran on top of it. Nothing in the substrate produced that number, and nothing would have noticed if it had been a different number. The mass was a free dial turned by hand.

On the depth rubric used throughout this paper that is an \emph{L0 hole}, a surface-level gap where a quantity that ought to emerge is instead put in. It is the same kind of gap the $g$-factor carried before it was measured from the substrate's own Dirac structure. A real substrate must produce the mass \emph{mechanism}. It must say what the mass is made of, why it can be nonzero at all, and what makes it vanish. The value can stay open. The mechanism cannot.

\subsection{The two chiralities}

Recall the local picture. A dock is the geometric cell, a here-now, and out of each dock run $24$ directed sites. Those $24$ sites are the $D_4$ root system, and under triality they split into three eights, written $8v + 8s + 8c$. The $8v$ piece is the vector, the photon sector. The two pieces $8s$ and $8c$ are the genuine half-integer-spin spinors. These two spinor sectors are the two \emph{chiralities} of the fermion, its two handednesses.

\begin{table}[ht]
\centering
\begin{tabular}{@{}ll@{}}
\toprule
sector & chirality \\
\midrule
$8s$ & one handedness \\
$8c$ & the other handedness \\
\bottomrule
\end{tabular}
\caption{The two spinor sectors of the dock are the two chiralities of the fermion.}
\label{tab:chiralities}
\end{table}

The Dirac mass term is a coupling between these two sectors. With no coupling the two chiralities are independent. Each one propagates on its own and never speaks to the other, and a fermion built from two independent chiralities is massless. With coupling the two sectors mix, each one feeds into the other, and that mixing \emph{is} the mass. The mass is not a thing added beside the chiralities. It is the bridge between them.

\begin{definition}[The chirality coupling]
The Dirac mass term is the only part of the dock's Hamiltonian that links the $8s$ sector to the $8c$ sector. When that link is present the two chiralities trade into one another at a definite rate. That rate is the mass.
\end{definition}

\begin{principle}[Handedness and the bridge]
Chirality is handedness. Mass is the bridge between the left hand and the right hand. With no bridge the two hands run free and massless.
\end{principle}

\subsection{Two measurements that must agree}

A mechanism is only convincing if the same number falls out of two unrelated definitions. The mass is read off the substrate in two independent ways. If the mechanism is real they are forced to give the same value. If they ever disagreed, the picture of mass as a chirality coupling would be wrong.

\subsubsection{The dispersion mass}

Take the Dirac operator $H$ on the dock and square it. The squared operator is the relativistic scalar at every momentum,
\[
H^2 = p^2 + m^2 .
\]
This is the same hinge that produces the Lorentz structure elsewhere in the paper. The point worth stressing is that $H$ squaring to a clean scalar is \emph{emergent}. It is not assumed. The directional law of the dock, collide then stream, is a definite discrete object, and when its Dirac operator is squared the relativistic form appears on its own. The $m$ that sits in that formula is read straight off the spectrum. Call it the \emph{dispersion mass}.

\subsubsection{The chirality-coupling mass}

Now take the commutator of the Hamiltonian with $\gamma_5$, the chirality operator. The operator $\gamma_5$ measures handedness. Everything in $H$ commutes with $\gamma_5$ except the mass term, which is the single piece that mixes the two handednesses. So the size of the commutator $[H, \gamma_5]$ measures the chirality coupling directly, and nothing else.

\begin{principle}[The commutator is the mass]
Everything in $H$ commutes with $\gamma_5$ except the mass. So the failure to commute is exactly the mass. The size of $[H, \gamma_5]$ is the chirality-coupling mass.
\end{principle}

\subsubsection{The agreement}

\begin{table}[ht]
\centering
\begin{tabular}{@{}lll@{}}
\toprule
measurement & what it reads & result \\
\midrule
dispersion mass & $H^2 = p^2 + m^2$ & $m$ \\
chirality-coupling mass & size of $[H, \gamma_5]$ & $m$ \\
\bottomrule
\end{tabular}
\caption{Two independent ways to read the mass off the substrate, and the single number both return.}
\label{tab:two-masses}
\end{table}

The two measurements agree to machine precision. One reads the mass off the dispersion, the slope of the spectrum at every momentum. The other reads it off the chirality coupling, the size of the one commutator that fails. These definitions share no algebra. They could have disagreed by any amount. They match exactly.

\begin{result}[Two roads, one number]
The dispersion mass from $H^2 = p^2 + m^2$ and the chirality-coupling mass from the size of $[H, \gamma_5]$ agree to machine precision \cite{pollard2026vibetest}. A coincidence would not survive two independent definitions matching that closely. The mass is the same whether read off the dispersion or off the chirality coupling.
\end{result}

That agreement is the heart of the result. It is what raises this from a slot-filling exercise to a genuine measurement.

\subsection{The control, the Weyl limit}

A measurement needs a control, a knob you can turn that could break the claim. Here the control is to switch the coupling off.

At zero coupling the two chiralities decouple. The $8s$ and the $8c$ no longer talk to one another. Each runs as an independent massless mode, and the fermion becomes a massless Weyl fermion, a particle of a single handedness moving at the speed of light.

\begin{table}[ht]
\centering
\begin{tabular}{@{}lll@{}}
\toprule
coupling & chiralities & mass \\
\midrule
on & $8s$ and $8c$ mixed & $m$, both measures agree \\
off & $8s$ and $8c$ independent & exactly zero, both measures \\
\bottomrule
\end{tabular}
\caption{The Weyl limit. Turning the coupling off sends both mass measurements to zero together.}
\label{tab:weyl-limit}
\end{table}

The decisive fact is that both meters fall to zero \emph{together}. The dispersion mass and the chirality-coupling mass vanish at the same point, the point of zero coupling. This is the control that could have failed. If only one of the two measures went to zero, the two would not have been measuring the same thing, and the chirality-coupling picture would collapse. They both vanish, and exactly together.

\begin{result}[The Weyl limit confirms the mechanism]
At zero coupling the $8s$ and $8c$ sectors decouple and a massless Weyl fermion is left. Both the dispersion mass and the chirality-coupling mass vanish there, on both meters at once. That joint vanishing is the proof that the mass \emph{is} the coupling, not a separate energy that merely tracks it.
\end{result}

\subsection{Why this counts as a real measurement}

The depth rubric reserves its L2 level for a genuine measurement equipped with a discriminating control. This result earns that level on four counts.

\begin{itemize}
\item The mass is read off the substrate, never assigned.
\item It is measured two independent ways that agree to machine precision.
\item The relativistic form $H^2 = p^2 + m^2$ is emergent, not assumed.
\item There is a clean control, the Weyl limit, where both measures vanish together.
\end{itemize}

Through all of it the base stays discrete and deterministic. The dock's $24$ sites carry ternary tones, the knit collides and streams them by a fixed reversible table, and time ticks one beat at a time. The Dirac operator and its $\gamma_5$ commutator are definite objects built on that definite mesh. Nothing continuous and nothing random is smuggled in. The mass falls out of the discrete substrate.

\subsection{Mechanism, not value}

One boundary must be stated plainly. This derives the mass \emph{mechanism}. It does not derive the mass \emph{value}.

The value is the Yukawa parameter, and it stays open. The substrate fixes how a mass can arise, the $8s$-$8c$ coupling carrying the relativistic form, but it does not fix how large that coupling is. The same boundary holds for the gauge coupling. Scanning the coupling shows the bare knit treats it as a free overall scale with no distinguished value anywhere, vanishing only at zero. A pure scale fixes a form, not a magnitude.

\begin{principle}[Form fixed, value free]
The knit fixes how mass arises, the $8s$-$8c$ coupling with the relativistic form. It does not fix how big the mass is. The magnitude is a free input the bare law does not contain.
\end{principle}

This mirrors the $g$-factor exactly. There the mechanism is derived and measured while the value boundary holds at the same place. Mechanism earned, value left to a principle the bare knit does not supply.

\subsection{Where the generation story sits}

A natural next question follows at once. Triality permutes the three eights. Do those three slots give three different masses, the three observed generations of matter? The honest answer is no, not yet, and the reason is worth setting down carefully because it is a place where the substrate gives a real structure but stops short of the full claim.

The naive reading does not work. The $8v$, $8s$, $8c$ triple is a vector plus two chiralities of \emph{one} generation, not three generations. Triality relating the three eights is not the same as three copies of a fermion.

The deeper path runs through the substrate's exceptional symmetry, and it is computed rather than asserted. The $24$ directions are the long roots of $F_4$, and $F_4$ is the automorphism group of the exceptional Jordan algebra $J_3(\mathbb{O})$, the $27$-dimensional Albert algebra. Its rank-three frame is computed directly, three primitive orthogonal idempotents summing to the identity. The three is \emph{forced}. The Jordan identity holds for Hermitian octonion matrices at size three and below, but fails at size four because the octonions are non-associative. The size-four failure is the control. The substrate's own symmetry forces an exactly three-fold exceptional structure.

\begin{proposition}[A forced three-fold structure]
The $24$ directions of the dock are the long roots of $F_4$, the automorphism group of the exceptional Jordan algebra $J_3(\mathbb{O})$. Its rank-three frame is three orthogonal idempotents summing to the identity, and the rank three is forced by the failure of the Jordan identity at size four, a consequence of octonion non-associativity.
\end{proposition}

The next necessary ingredient is a family symmetry relating the three slots, and it exists and is exact. The symmetric group $S_3$ permuting the slots acts by genuine Jordan automorphisms, with a cyclic order-three generation triality. But the same computation exposes the honest gap. Because the $S_3$ is exact, the three slots are \emph{degenerate}. Every quantity the algebra supplies is equal across them. The substrate gives a family symmetry but no family breaking, and three identical slots are not three different masses.

\begin{principle}[Symmetry without breaking]
The substrate gives a family symmetry but no family breaking. Three identical slots are not three masses. Splitting them is Boyle's open conjecture, and it stays open.
\end{principle}

The mass mechanism in this section is solid. The split of the degenerate triple into three distinct masses is the honest residual, a problem open in physics and not merely in the present code.

\subsection{Why it matters}

Mass is the second thing a matter substrate must earn rather than assume. The first was spin, and the chain now runs cleanly across three sections.

\begin{itemize}
\item The dock's sites carry two chiralities, the $8s$ and the $8c$.
\item Those chiralities propagate relativistically, with $H^2 = p^2 + m^2$ emerging from the discrete law.
\item Their coupling is the mass, measured two agreeing ways, vanishing into a Weyl fermion when off.
\end{itemize}

\begin{result}[Mass earned on the substrate]
On $\{3,4,3,4\}$ the fermion mass is the coupling of the two chiralities $8s$ and $8c$. It is read off the substrate by dispersion and by the $\gamma_5$ commutator, the two agree to machine precision, and both vanish together in the Weyl limit. The mechanism is derived and measured. The value is the open Yukawa parameter, and the split into three generations is the open residual.
\end{result}

\begin{principle}[The marriage of the two hands]
The mass is the marriage of the two hands of the fire. Unmarry them and the fire runs massless at the speed of light.
\end{principle}
